\section{\uppercase{Related Work}} \label{sec2}
Research on adaptive RAG, which dynamically selects retrieval methods and hyperparameters according to query characteristics, has grown rapidly~\cite{Asai23,fu24,Tang25,Rishi25,Jeong24}. Generally, high-accuracy techniques such as graph-based RAG are effective for complex queries, whereas applying complex pipelines to simple queries incurs unnecessary computational cost. Conversely, processing complex queries with simplistic approaches (e.g., NaiveRAG) or without retrieval leads to degraded response quality~\cite{Jeong24}. Adaptive RAG aims to improve response quality while optimizing cost by dynamically selecting retrieval strategies and hyperparameters based on query complexity.

Most existing studies rely on LLM self-judgment or online reinforcement learning. Self-RAG~\cite{Asai23} enables LLMs to self-evaluate and control retrieval during generation. AutoRAG-HP~\cite{fu24} optimizes multiple RAG hyperparameters using online reinforcement learning. Tang et al.~\cite{Tang25} proposed an automated retrieval selection method for graph-based RAG based on query properties. HyPA-RAG~\cite{Rishi25} dynamically optimizes multiple RAG hyperparameters according to query characteristics. Furthermore, various methods for adaptively optimizing RAG processing workflows and chunk combinations have been proposed. For example, CARROT~\cite{wang2026carrot} employs online Monte Carlo Tree Search (MCTS) to optimize retrieved chunk combinations and ordering under dynamic cost constraints. uRAG~\cite{salemi2024urag} provides a unified retrieval and re-ranking engine shared across diverse downstream RAG systems, optimizing document ranking based on agent feedback. RAGRouter~\cite{zhang2025ragrouter} models whether an LLM can effectively process retrieved document context and dynamically routes queries among multiple generator LLMs with varying capabilities and costs. AdaGCRAG~\cite{zhang2025adagcrag} controls the scope and depth of adaptive entity-relation extraction in graph-based RAG to improve efficiency in lightweight LLM environments.

In contrast to these prior works, our proposed method unifiedly ranks a complete configuration space of 96 concrete candidate configurations, spanning RAG paradigms (Graph / Non-Graph), chunk sizes, chunk overlaps, and combinations of generator and embedder models. Furthermore, because our method relies on offline learning that requires no additional LLM invocations or online tree search (e.g., MCTS) during inference, decision latency is extremely low. Crucially, by simply adjusting priority weights $W$, the system can flexibly handle multi-objective trade-offs (precision, recall, F1 score, and execution latency).

